% -*-latex-*-
% RMIT requirement: technical abstract, <200 words (counted as printed).

Large language models (LLMs) increasingly score student essays, yet
models trained mainly on Western text risk treating Anglo-American
rhetoric as the standard and penalising other valid traditions. This thesis formalises the mechanism as \emph{evaluative
epistemic filtering} and develops CAMPS (Culturally-Aware
Multi-Perspective Scoring), an inference-time framework where
agents calibrated to distinct cultural readings of the rubric score
each essay independently, a Bias Divergence Score (BDS) measuring
disagreement. Across 200 topic-controlled
ICNALE essays and four model families, every baseline model scored
native-English above upper-band Southeast Asian writers
(0.19--1.18 bands; all 95\% intervals exclude zero). As the cohorts
were not matched on judged quality, this is an unadjusted disparity,
not an isolated cultural bias. A
three-agent panel reduced the disparity by 23--28\% on the two most
disparate models, by marking both cohorts down (natives more) and by no more than
its Southeast Asian lens alone; it amplified the disparity on the least disparate model. Five agents
added nothing on GPT-4o and widened the gap on Claude. Reweighting showed no trade-off with band association. BDS missed its detection target (held-out sensitivity 0.06 and
0.00 at 5\% false-positive rate). Culturally framed prompts change LLM scoring
in model-dependent ways; cohort gaps and agent disagreement are
imperfect fairness proxies.
